% Part: turing-machines
% Chapter: undecidability
% Section: representing-tms

\documentclass[../../../include/open-logic-section]{subfiles}

\begin{document}

\olfileid{tur}{und}{rep}
\olsection{ટ્યુરિંગ મશીનોનું નિરૂપણ}

\begin{explain}
ટ્યુરિંગ મશીનો અને તેમનું વર્તન પ્રથમ-ક્રમના તર્કના
!!a{sentence} વડે નિરૂપવા માટે યોગ્ય ભાષા વ્યાખ્યાયિત
કરવી પડશે. આ ભાષાના બે ભાગ છે: મશીનનાં સંરૂપણોનું
વર્ણન કરવા માટેનાં !!{predicate}s, અને સંગણનનાં પગલાં
(``ક્ષણો'') તથા ટેપ પરનાં સ્થાનોને અંક આપવા માટેની
અભિવ્યક્તિઓ.

બે જાતનાં !!{predicate}s લઈએ છીએ, અને બંને દ્વિસ્થાનિક
છે: દરેક અવસ્થા~$q$ માટે !!a{predicate}~$\Obj Q_q$,
અને દરેક ટેપ-ચિહ્ન~$\sigma$ માટે
!!a{predicate}~$\Obj S_\sigma$. પહેલાંના વિધાનો
આપણને~$M$ની અવસ્થા અને તેના ટેપ-હેડનું સ્થાન
વર્ણવવા દે છે; બીજાં વિધાનો ટેપની સામગ્રીનું વર્ણન
કરવા દે છે.

ટેપ-હેડનાં સ્થાનો અને થઈ ચૂકેલાં પગલાંની સંખ્યા
વ્યક્ત કરવા સંખ્યાઓ વ્યક્ત કરવાની રીત જોઈએ. તે માટે
!!a{constant}~$\Obj 0$ અને $1$-સ્થાનિક અનુગામી
વિધેય~$\prime$ વાપરીએ છીએ. રૂઢિ મુજબ તેને તેના
આર્ગ્યુમેન્ટની \emph{પછી} લખીએ છીએ (અને કૌંસ
છોડી દઈએ છીએ).

દરેક સંખ્યા $n$ માટે પ્રમાણભૂત પદ~$\num{n}$ છે,
એટલે~$n$નું \emph{આંકિક પદ}, જે~$\Lang L_M$માં
તેનું નિરૂપણ કરે છે. $\num{0}$ એ $\Obj 0$,
$\num{1}$ એ $\Obj 0'$, $\num{2}$ એ $\Obj 0''$,
વગેરે છે. વધુ ઔપચારિક રીતે:
\begin{align*}
\num{0} & = \Obj 0 \\
\num{n+1} &= \num{n}'
\end{align*}
પદ $\num{0}$, એટલે $\Obj 0$, ટેપના સૌથી ડાબા
સ્થાનને તેમજ પ્રથમ પગલા પહેલાંના સમયને (પ્રારંભિક
સંરૂપણને) નામ આપે છે. પદ $\num{1}$, એટલે
$\Obj 0'$, સૌથી ડાબા ખાનાની જમણી બાજુના ખાનાને
અને પ્રથમ પગલા પછીના સમયને નામ આપે છે; આગળ
પણ એમ જ છે.

ટેપનાં સ્થાનોનો ક્રમ (અર્થ ``ની ડાબે'') તથા
સંગણનનાં પગલાંનો ક્રમ (અર્થ ``પહેલાં'') બંને વ્યક્ત
કરવા !!a{predicate}~$<$ પણ લઈએ છીએ.

ભાષા નક્કી થયા પછી~$!T(M, w)$નાં ``સ્વયંસિદ્ધ વાક્યો''
યાદીબદ્ધ કરીએ: એટલે એવા !!{sentence}s, જે સાથે
મળીને ઇનપુટ~$w$ પર ચાલતાં~$M$નું વર્તન વર્ણવે.
$\Obj 0$, $\prime$ અને $<$ પર શરતો મૂકતાં !!{sentence}s,
પ્રારંભિક સંરૂપણ વર્ણવતાં !!{sentence}s અને ચોક્કસ સૂચના
ચલાવ્યા પછી~$M$નું સંરૂપણ કેવું હોય તે વર્ણવતાં
!!{sentence}s હશે.
\end{explain}

\begin{defn}
  \ollabel{defn:tm-descr}
ટ્યુરિંગ મશીન $M = \tuple{Q, \Sigma, q_0, \delta}$
આપેલું હોય, તો ભાષા~$\Lang L_M$માં નીચેનું આવે છે:
\begin{enumerate}
\item દરેક અવસ્થા~$q \in Q$ માટે દ્વિસ્થાનિક
  !!{predicate} $\Obj Q_q(x, y)$. સહજ રીતે,
  $\Obj Q_q(\num{m}, \num{n})$ વ્યક્ત કરે છે કે
  ``$n$ પગલાં પછી~$M$,~$m$મું ખાનું વાંચતાં
  અવસ્થા~$q$માં છે.''
\item દરેક ચિહ્ન~$\sigma\in \Sigma$ માટે દ્વિસ્થાનિક
  !!{predicate} $\Obj S_\sigma(x, y)$. સહજ રીતે,
  $\Obj S_\sigma(\num{m}, \num{n})$ વ્યક્ત કરે છે કે
  ``$n$ પગલાં પછી~$m$મા ખાનામાં ચિહ્ન~$\sigma$ છે.''
\item એક !!{constant} $\Obj 0$
\item એક એકસ્થાનિક !!{function} $\prime$
\item એક દ્વિસ્થાનિક !!{predicate} $<$
\end{enumerate}
\end{defn}

ટ્યુરિંગ મશીન~$M$, ઇનપુટ
$w = \sigma_{i_1}\dots\sigma_{i_k}$ પર કેવી રીતે
કામ કરે તે વર્ણવતાં !!{sentence}s નીચે મુજબ છે:
\begin{enumerate}
\item સંખ્યાઓ અને~$<$નું વર્ણન કરતાં સ્વયંસિદ્ધ વાક્યો:
\begin{enumerate}
\item દરેક સંખ્યા તેના અનુગામી કરતાં નાની છે એવું
કહેતું !!^a{sentence}:
\[
\lforall[x][x < x']
\]
\item $<$ સંક્રામક છે તેની ખાતરી આપતું !!^a{sentence}:
\[
\lforall[x][\lforall[y][\lforall[z][
      ((x < y \land y < z) \lif x < z)]]]
\]
\end{enumerate}
\item ઇનપુટ-સંરૂપણનું વર્ણન કરતાં સ્વયંસિદ્ધ વાક્યો:
\begin{enumerate}
\item મશીન શરૂ થાય તે પહેલાં, એટલે $0$~પગલાં પછી,
  $M$ પ્રારંભિક અવસ્થા~$q_0$માં છે અને ખાનું~$1$
  વાંચે છે:
\[
\Obj Q_{q_0}(\num{1}, \num{0})
\]
\item પ્રથમ $k+1$ ખાનાંમાં ચિહ્નો $\TMendtape$,
  $\sigma_{i_1}$, \dots, $\sigma_{i_k}$ છે:
\[
\Obj S_\TMendtape(\num{0}, \num{0}) \land
\Obj S_{\sigma_{i_1}}(\num{1}, \num{0}) \land
\dots \land
\Obj S_{\sigma_{i_k}}(\num{k}, \num{0})
\]
\item બાકીની ટેપ ખાલી છે:
\[
\lforall[x][(\num{k} < x \lif \Obj S_\TMblank(x, \num{0}))]
\]
\end{enumerate}
\item એક સંરૂપણમાંથી આગળના સંરૂપણમાં થતું સંક્રમણ
  વર્ણવતાં સ્વયંસિદ્ધ વાક્યો:

આગળ માટે, $!A(x, y)$ એ નીચેના સ્વરૂપનાં બધાં
!!{sentence}sનું સંયોજન માનો:
\[
\lforall[z][
  (((z < x \lor x < z) \land \Obj S_\sigma(z, y))
  \lif \Obj S_\sigma(z, y'))]
\]
અહીં $\sigma \in \Sigma$. $!A(\num{m},\num{n})$
વડે આપણે કહીએ છીએ કે ``$m$મા ખાના સિવાય,
$n+1$ પગલાં પછીની ટેપ $n$ પગલાં પછી જેવી જ છે.''
\begin{enumerate}
\item \ollabel{rep-right} દરેક સૂચના
  $\delta(q_i, \sigma) =
  \tuple{q_j, \sigma', \TMright}$ માટે !!{sentence}:
\begin{align*}
& \lforall[x][\lforall[y][(
   (\Obj Q_{q_i}(x, y) \land \Obj S_{\sigma}(x, y)) \lif {}]] \\
&\qquad   (\Obj Q_{q_j}(x', y') \land \Obj S_{\sigma'}(x, y') \land
!A(x, y)))
\end{align*}
એનો અર્થ એ છે કે~$y$ પગલાં પછી મશીન
અવસ્થા~$q_i$માં, ચિહ્ન~$\sigma$ ધરાવતું ખાનું~$x$
વાંચે છે, તો $y+1$ પગલાં પછી તે ખાનું~$x+1$ વાંચે છે,
અવસ્થા~$q_j$માં છે, ખાનું~$x$ હવે~$\sigma'$ ધરાવે છે,
અને~$x$ સિવાયનું દરેક ખાનું~$y$ પગલાં પછી જે ચિહ્ન
ધરાવતું હતું તે જ ધરાવે છે.

\item \ollabel{rep-left} દરેક સૂચના
  $\delta(q_i, \sigma) =
  \tuple{q_j, \sigma', \TMleft}$ માટે !!{sentence}:
\begin{align*}
& \lforall[x][\lforall[y][
    ((\Obj Q_{q_i}(x', y) \land \Obj S_{\sigma}(x', y)) \lif {}]]\\
& \qquad   (\Obj Q_{q_j}(x, y') \land \Obj S_{\sigma'}(x', y') \land
!A(x', y))) \land {}\\
& \lforall[y][((\Obj Q_{q_i}(\num{0}, y) \land \Obj S_{\sigma}(\num{0},
    y)) \lif {}]\\
& \qquad (\Obj Q_{q_j}(\num{0}, y') \land \Obj S_{\sigma'}(\num{0},
  y') \land !A(\num{0}, y)))
\end{align*}
\sourcecorrection{OLUN-002}{સ્થિર અંગ્રેજી મૂળની
ડાબે-ચાલ સૂચનાની પહેલી શાખામાં લખાતું ખાનું
$x'$ છે, પરંતુ તેનું $!A(x,y)$ એ ખાનાને
યથાવત્ રહેવાની શરતમાંથી બાકાત રાખતું નથી.
અહીં લખાતા ખાનાને બાકાત રાખવા $!A(x',y)$ કર્યું છે.}
આ કેવી રીતે કામ કરે છે તે વિચારો: હવે આપણે
``ખાનું~$x$ વાંચતું હોય તો \dots''થી શરૂ કરતા નથી,
પણ ``ખાનું $x+1$ વાંચતું હોય તો \dots''થી શરૂ કરીએ
છીએ. ડાબે જવાથી આગળના પગલે મશીન ખાનું~$x$
વાંચે છે. પરંતુ લખાણ ખાનું~$x+1$ પર થાય છે.
બાદબાકી કે પૂર્વગામી વિધેય ન હોવાથી આ રીતે કરીએ
છીએ.

નોંધો કે $x+1$ સ્વરૂપની સંખ્યાઓ $1$, $2$, \dots છે;
એટલે મશીન ખાનું~$0$ વાંચતું હોય અને ડાબે જવાનું હોય
તે કિસ્સો અહીં આવતો નથી (અને તે ડાબે જઈ પણ ન શકે,
એટલે એ જ સ્થાને રહે છે). બીજું સંયોજન એ ખાસ કિસ્સો
આવરે છે: $y$ પગલાં પછી મશીન અવસ્થા~$q_i$માં
ખાનું~$0$ વાંચે અને ખાનું~$0$માં ચિહ્ન~$\sigma$
હોય, તો $y+1$ પગલાં પછી પણ તે ખાનું~$0$ વાંચે છે,
હવે અવસ્થા~$q_j$માં છે, ખાનું~$0$માં ચિહ્ન~$\sigma'$
છે, અને ખાનું~$0$ સિવાયનાં ખાનાંમાં~$y$ પગલાં
પછી હતાં તે જ ચિહ્નો રહે છે.
\item \ollabel{rep-stay} દરેક સૂચના
  $\delta(q_i, \sigma) =
  \tuple{q_j, \sigma', \TMstay}$ માટે !!{sentence}:
\begin{align*}
& \lforall[x][\lforall[y][(
   (\Obj Q_{q_i}(x, y) \land \Obj S_{\sigma}(x, y)) \lif {}]] \\
&\qquad   (\Obj Q_{q_j}(x, y') \land \Obj S_{\sigma'}(x, y') \land
!A(x, y)))
\end{align*}
\end{enumerate}
\end{enumerate}
ટ્યુરિંગ મશીન~$M$ અને ઇનપુટ~$w$ માટે ઉપરનાં
બધાં !!{sentence}sનું સંયોજન~$!T(M, w)$ માનો.

~$M$ છેવટે અટકે છે એમ વ્યક્ત કરવા એવું !!{sentence}
જોઈએ જે કહે કે ``કેટલાંક પગલાં પછી સંક્રમણ વિધેય
અવ્યાખ્યાયિત હશે.'' $X$ એ એવાં બધા
$\tuple{q, \sigma}$ જોડોનો ગણ માનો કે જેમના માટે
~$\delta(q, \sigma)$ અવ્યાખ્યાયિત છે. હવે~$!E(M, w)$
એ નીચેનું !!{sentence} માનો:
\[
\lexists[x][\lexists[y][(\bigvee_{\tuple{q, \sigma} \in
      X}(\Obj Q_q(x, y) \land \Obj S_\sigma(x, y)))]]
\]

ખાસ અટકવાની અવસ્થા~$h$ ધરાવતું ટ્યુરિંગ મશીન
વાપરીએ, તો વાત વધુ સરળ બને છે: તે વખતે
!!{sentence}~$!E(M, w)$
\[
\lexists[x][\lexists[y][\Obj Q_h(x, y)]]
\]
મશીન છેવટે અટકે છે એમ વ્યક્ત કરે છે.

\begin{prop}
\ollabel{prop:mlessk}
જો $m < k$, તો $!T(M, w) \Entails \num{m} < \num{k}$.
\end{prop}

\begin{proof}
કસરત.
\end{proof}

\begin{prob}
\olref[tur][und][rep]{prop:mlessk} સાબિત કરો.
(સૂચન: $k-m$ પર આગમન વાપરો.)
\end{prob}

\end{document}
